% Modernised reading — fonds Alexandre Grothendieck, shelfmark 41, pages 1-13.
%
% This is an INTERPRETATION, not a transcription. Everything the critical
% apparatus carried moves into footnotes. In any dispute of substance the
% transcription governs: whoever cites Grothendieck must cite batch-01.fr.tex.
%
% Pass: Opus 5 (claude-opus-5), 2026-09-13 — first pass, unchecked against the
% pages by a human. The folder was read whole; it holds one batch, pages 1-13,
% of which page 2 (the blank back of the title leaf) carries nothing.
%
% What the folder is. A draft section, under a title of his, on the étale
% Y-scheme M whose geometric points over y are the irreducible components of
% the geometric fibre of f : X -> Y: when it exists, when it is separated, when
% it is proper. The hand is drafting and most of the connective prose is
% illegible; the definitions, the statements and the displayed conditions are
% what carry, and the reading is built on them.
%
% Notation fixed for the whole document. f : X -> Y of finite type, Y locally
% noetherian where needed. X_y the fibre, X_{\bar y} a geometric fibre.
% Irr(Z) the set of irreducible components of a space Z. His « rel^t
% irréductible » is written « géométriquement irréductible » throughout, on the
% evidence of page 4, where a component is « rel^t irréd. » exactly when the
% point of M indexing it is rational over kappa(y); his « abs^t » (page 8) is
% left as « irréductible » and the ambiguity is footnoted. His E_{X/Y}(y) is
% kept. « Monogène » is his word and is kept, with its two definitions (page 3
% and pages 11-12) reconciled on the spine.
%
% Substantive departures from the pages, each footnoted where it occurs:
% — page 4, « Spec(Omega_y) »: the reduction it serves is to a local base, and
%   the reading takes the local scheme Spec(O_{Y,y}); over a geometric point
%   the statement would be empty.
% — page 4, descent along the completion: effective for M separated (quasi-
%   finite separated => quasi-affine), not guaranteed in schemes for M
%   non-separated; the leaf does not distinguish.
% — page 9, N.B.: the injection runs from the components of the special fibre
%   into those of the generic fibre, through the sections of M over a strictly
%   henselian trait; the reading states it so.
% — page 9, « un ouvert étale »: what condition (iii) gives is the existence
%   half of the valuative criterion, hence M proper, hence finite étale; the
%   word on the leaf does not carry that sense and is footnoted.
% — page 11, Prop. 1: (i) <=> (ii) holds with no hypothesis (phi^{-1}(U) =
%   f(U), his own formula); (ii) <=> (iii) needs f locally of finite
%   presentation. The reading says which implication uses what.
% — page 13, the struck Lemme: as legible, it fails for a flat f without a
%   further hypothesis; a counterexample is given in a footnote.
%
% The three valuative conditions on the spine (uniqueness of the generisation,
% separation, existence) are this reading's analysis of why the Théorème has
% the shape it has. They are ours and are marked as such; the leaves' own
% conditions are largely illegible, and the reading does not claim that the
% Théorème, as the leaf states it, is proved or even fully stated.
%
% What the folder announces and never establishes: the proof of the Théorème
% of page 8; the Lemme of page 7 beyond its conclusion; conditions a) b) c) of
% page 10; the representability asked on page 12 (« Est-il représentable ??? »);
% the two conclusions of page 13, closed by « Vérifier ! ».
\documentclass[11pt,a4paper]{article}
% Le préambule commun à toutes les transcriptions du fonds Grothendieck.
%
% Il tient en une idée : **l'incertitude se compose, elle ne se résout pas.**
% Une transcription qui lisse un mot douteux a détruit la seule chose pour
% laquelle on la faisait — un lecteur doit pouvoir savoir, à chaque endroit,
% ce qui était sur le feuillet et ce qui a été ajouté. Les six macros qui
% suivent sont donc l'appareil critique, et non de la décoration.
%
% Les mêmes commandes sont reconnues par `scripts/render.mjs`, qui produit la
% vue de lecture du site. Une macro ajoutée ici doit l'être aussi là-bas,
% faute de quoi le rendu échouera bruyamment — ce qui est voulu : un rendu
% silencieusement faux est pire qu'un rendu absent.

\ProvidesPackage{grothendieck}[2026/08/08 Transcriptions du fonds Grothendieck]

\usepackage{amsmath,amssymb,amsthm}
\usepackage{mathrsfs}
\usepackage{stmaryrd}
% Les diagrammes commutatifs sont partout dans ce fonds : c'est la langue
% même de la géométrie algébrique de Grothendieck, et une transcription qui
% les rendrait en prose ne transcrirait rien.
\usepackage{tikz-cd}
% Français par défaut — c'est la langue des feuillets — mais l'anglais est
% chargé aussi : la traduction et le résumé sont des documents anglais, et la
% typographie française (espaces avant les deux-points) y serait une faute.
% Ces documents basculent par \selectlanguage{english} en tête de corps.
\usepackage[english,french]{babel}
\usepackage{fontspec}
\usepackage{geometry}
\geometry{a4paper,margin=25mm,marginparwidth=28mm}
\usepackage{marginnote}
\usepackage{xcolor}
% ulem, not soul. `soul`'s \st and \ul are built on a character-by-character
% scan that XeLaTeX's font machinery breaks: the document fails with an
% undefined control sequence rather than degrading. ulem does the same two jobs
% and is Unicode-engine safe. `normalem` stops it hijacking \emph, which the
% transcriptions use for Grothendieck's own emphasis.
\usepackage[normalem]{ulem}
\usepackage{hyperref}
\hypersetup{colorlinks=true,linkcolor=black,urlcolor=black,pdfborder={0 0 0}}

% --- Métadonnées du lot -----------------------------------------------------
% Reprises telles quelles par le script de rendu, qui les affiche en tête de
% la vue de lecture. Elles fixent ce que le fichier prétend couvrir : sans
% elles, une transcription flotte sans point d'attache dans un fonds de seize
% mille feuillets.

\newcommand{\folder}[1]{\gdef\grothFolder{#1}}
\newcommand{\batch}[1]{\gdef\grothBatch{#1}}
\newcommand{\leaves}[2]{\gdef\grothFirst{#1}\gdef\grothLast{#2}}
\newcommand{\foldertitle}[1]{\gdef\grothFoldertitle{#1}}
\gdef\grothFolder{?}\gdef\grothBatch{?}\gdef\grothFirst{?}\gdef\grothLast{?}\gdef\grothFoldertitle{}

% --- Le résumé --------------------------------------------------------------
%
% Il ouvre la lecture modernisée, sous le titre « Résumé », et il est composé
% en petit. Ce n'est pas un ornement : c'est le seul passage du document qui
% ne soit pas des mathématiques, et un lecteur doit pouvoir le reconnaître
% avant de l'avoir lu — puis le sauter s'il connaît déjà le sujet, ou s'y
% arrêter s'il ne le connaît pas. Le filet qui le suit marque la reprise du
% corps.

\newenvironment{resume}
  {\small\setlength{\parskip}{0.45em}}
  {\par\medskip\hrule\medskip}

% --- Mots-clés ---------------------------------------------------------------
%
% En anglais, en clôture du résumé de la lecture modernisée : le vocabulaire
% moderne sous lequel chercher ce que les feuillets construisent. Le manifeste
% du site (`npm run manifest`) les extrait pour étiqueter la cote entière —
% cette ligne est la source unique des tags, il n'y a pas de fichier à côté.
\newcommand{\keywords}[1]{\par\smallskip\noindent
  {\footnotesize\textsc{Keywords} --- \textit{#1}}\par}

% --- L'appareil critique ----------------------------------------------------

% Le numéro de feuillet, en marge. C'est l'ancre qui permet de régler un
% désaccord sur une formule en regardant *un* feuillet plutôt qu'une chemise
% de six cents. Le site s'en sert aussi pour tourner les pages du fac-similé
% quand on fait défiler la transcription.
\newcommand{\leaf}[1]{%
  \marginnote{\footnotesize\color{gray!70}\textsf{#1}}%
  \label{leaf:#1}\ignorespaces}

% Illisible : on ne devine pas. Un mot inventé qui se lit comme les autres est
% le pire résultat possible.
\newcommand{\ill}{\textcolor{red!60!black}{[\,\ldots\,]}}

% Lecture douteuse : le mot est donné, et signalé comme incertain.
\newcommand{\uncertain}[1]{\uline{#1}}

% Ajout de l'éditeur — une lettre manquante, un mot sous-entendu.
\newcommand{\add}[1]{\textcolor{blue!50!black}{[#1]}}

% Biffé par Grothendieck lui-même. Ce qu'il a rayé fait partie du document :
% une rature dit souvent où la pensée a changé de direction.
\newcommand{\struck}[1]{\sout{#1}}

% Note du transcripteur, sur l'état matériel du feuillet.
\newcommand{\note}[1]{\par\smallskip{\small\color{gray!60!black}\textit{[#1]}}\par\smallskip}

% Note marginale de Grothendieck, distincte du corps du texte.
\newcommand{\marginal}[1]{\par\smallskip\noindent\hspace{1em}%
  {\small\color{gray!60!black}#1}\par\smallskip}

% --- Titre ------------------------------------------------------------------

\AtBeginDocument{%
  \begin{center}
    {\large\bfseries Fonds Alexandre Grothendieck --- cote n\textsuperscript{o}~\grothFolder}\\[2pt]
    {\itshape\grothFoldertitle}\\[4pt]
    {\small feuillets \grothFirst--\grothLast\ (lot \grothBatch)}\\[2pt]
    {\footnotesize\color{gray!60!black}Transcription établie d'après le fac-similé numérisé
     par l'Université de Montpellier.\\
     Les lectures incertaines sont soulignées, les illisibles notées [\,\ldots\,],
     les ajouts entre crochets.}
  \end{center}
  \vspace{1em}}

\endinput


\folder{41}
\pages{1}{13}
\dating{[années 1960-1970]}
\watermark{Édition de démonstration\\interprétation personnelle de l'œuvre}
\foldertitle{Construction de morphismes étales (Chap. V) : notes manuscrites (s.d.) — lecture modernisée du dossier entier}

\begin{document}

\section*{Résumé}

\begin{resume}
Une famille d'objets géométriques qui dépend d'un paramètre peut changer de
forme quand le paramètre varie. Une conique $xy = t$ est d'un seul tenant pour
$t \neq 0$ et se casse en deux droites pour $t = 0$ ; inversement, deux
morceaux séparés peuvent venir se recoller. Ces feuillets posent une question
très simple sur ce phénomène : peut-on \emph{compter} les morceaux --- les
composantes irréductibles --- de chaque membre de la famille d'une manière qui
dépende elle-même bien du paramètre ?

Compter ne suffit pas ; on veut \emph{étiqueter}. L'idée est de fabriquer un
second objet $M$, posé au-dessus de l'espace des paramètres, dont les points
au-dessus d'une valeur $y$ du paramètre soient exactement les composantes du
membre de la famille correspondant à $y$. C'est un revêtement, au sens où les
topologues l'entendent : au-dessus de chaque point, un ensemble discret de
feuillets, qui se prolongent quand on bouge un peu. En géométrie algébrique ce
genre d'objet s'appelle un schéma \emph{étale} sur la base, et le titre du
dossier le dit : ce sont des morphismes étales associés aux composantes
irréductibles des fibres.

La difficulté est que les feuillets d'un tel revêtement ne se comportent pas
tous bien. Quand une composante se casse en deux au passage d'une valeur
spéciale, un feuillet se dédouble --- l'objet $M$ ressemble alors à une droite
dont on aurait doublé un point, et il n'est plus \emph{séparé}. Quand une
composante disparaît, un feuillet s'arrête net, et $M$ n'est plus
\emph{propre}. Le cœur du dossier (pages 8 à 10) est l'énoncé qui fait
correspondre ces pathologies de $M$ à la manière dont les composantes des
fibres se spécialisent les unes dans les autres : une composante générique ne
doit pas se casser en deux, deux composantes génériques ne doivent pas venir
se fondre, aucune ne doit s'évanouir. C'est la traduction, dans ce cadre, des
critères valuatifs par lesquels on teste la séparation et la propreté d'un
morphisme : on regarde ce qui se passe au-dessus d'un petit disque.

Le brouillon est difficile et souvent illisible, et il ne démontre pas son
théorème. Il en montre la genèse : une notion forgée pour l'occasion
(\emph{monogène}, pour une section qui choisit dans chaque fibre le point
générique d'une composante de manière continue), la question de sa
représentabilité, le passage à une base locale et complète, et l'essai d'une
construction, repris plusieurs fois. Le sujet s'appelle aujourd'hui l'espace
des composantes irréductibles d'un morphisme, et on le traite dans le cadre des
espaces algébriques, précisément parce que $M$ n'est pas toujours séparé.

\keywords{irreducible components of fibres, geometrically irreducible, universally open morphism, space of irreducible components, valuative criterion, strictly henselian trait}
\end{resume}

\section*{Le fil du dossier, et les conventions}

Le dossier est une seule section de rédaction, sous un titre de sa main, et il
avance en quatre temps.

\begin{enumerate}
\item[1.] \emph{La notion et la question} (page 3). Une section est
\emph{monogène} si elle choisit continûment le point générique d'une
composante géométriquement irréductible de chaque fibre ; ces sections forment
un foncteur de la base ; est-il représentable, et par quoi ?
\item[2.] \emph{Réduction et construction} (pages 4 à 7). Le problème est
local sur la base, se ramène à une base locale complète, et s'y construit
quand la fibre fermée est irréductible. Un Lemme résume ce qui a été obtenu.
\item[3.] \emph{Le théorème} (pages 8 à 10). Existence d'un $M$ étale séparé
sous deux conditions de spécialisation, puis d'un $M$ propre sous une
troisième, avec deux remarques.
\item[4.] \emph{Reprise} (pages 11 à 13). Retour au cas des fibres
irréductibles, où la continuité de la section générique équivaut à
l'ouverture du morphisme ; définitions arrêtées de « universellement
monogène » et de « monogène » ; la question de représentabilité reposée ;
un essai de lemme de spécialisation, barré.
\end{enumerate}

\emph{Conventions.} $f : X \to Y$ est de type fini. On note $X_y$ la fibre en
$y \in Y$, $X_{\bar{y}}$ une fibre géométrique, c'est-à-dire
$X \otimes_{Y} \Omega$ pour une extension algébriquement close $\Omega$ de
$\kappa(y)$, et $\mathrm{Irr}(Z)$ l'ensemble des composantes irréductibles d'un
espace $Z$. On garde sa notation $E_{X/Y}(y)$ pour l'ensemble des composantes
de la fibre $X_y$ qu'il qualifie de « relativement irréductibles », et l'on
écrit ici \emph{géométriquement irréductibles}.\footnote{Il écrit l'adverbe en
abrégé, « rel » suivi d'un t surélevé et souligné, et ne le définit nulle part.
Le sens retenu est fixé par la page 4 : la composante « rel.\ irréd. » qui
correspond à un point $\psi(y)$ de $M$ est celle pour laquelle $\psi(y)$ est
rationnel sur $\kappa(y)$, et un point de $M$ au-dessus de $y$ rationnel sur
$\kappa(y)$ indexe une composante que l'extension des scalaires ne casse pas.
La page 8 introduit l'abréviation parallèle « abs.\ » (absolument), insérée en
interligne au-dessus de « irréductibles » ; dans l'usage de Weil « absolument
irréductible » voulait précisément dire géométriquement irréductible, ce qui
ferait des deux adverbes des synonymes, et l'on ne comprendrait pas qu'il les
distingue. On lit donc ici « abs.\ irréductible » comme « irréductible »,
sans plus, et l'on signale que cette lecture est la nôtre.} « Monogène » est
son mot, et on le garde.

\emph{Deux définitions de « monogène ».} La page 3 appelle monogène une
application $\varphi : Y \to X$ telle qu'il existe un ouvert $U$ de $X$,
universellement ouvert sur $Y$, avec $\varphi(y)$ point générique de
$U \cap X_y$ pour tout $y$. La page 11 passe par un intermédiaire :
$U \to Y$ est \emph{universellement monogène} (fibres géométriquement
irréductibles, morphisme universellement ouvert), et $\varphi$ est monogène si
elle est la section des points génériques d'un tel $U$. Les deux disent la même
chose, et c'est la seconde qu'on adopte.

\emph{Ce que le dossier annonce et n'établit pas} : la démonstration du
théorème de la page 8 ; la plupart des hypothèses du Lemme de la page 7 ; les
conditions a), b), c) de la page 10 ; la représentabilité que la page 12
redemande avec trois points d'interrogation ; les deux conclusions de la page
13, qu'il fait suivre de « Vérifier ! ».

\pagerange{1}{1}
\section*{Le titre (page 1)}

Le premier feuillet ne porte que l'en-tête d'une rédaction : « Chap.\ VI »,
une ligne sur un paragraphe consacré à la descente et à ses applications, deux
mots biffés --- « Construction », « Étude » ---, et le titre, souligné sur deux
lignes : \emph{Morphismes étales associés à l'étude des composantes
irréductibles d'un morphisme ouvert}.\footnote{Le chiffre est un VI sur le
feuillet et non le V du titre de l'inventaire ; le premier des deux mots biffés
est celui que les archivistes ont retenu. Les deux mots qui suivent « Dans le
§ de la » ne sont pas assurés, « descente » étant une lecture incertaine et le
suivant illégible. Rien dans le dossier ne dit de quel ouvrage ce chapitre
devait faire partie ; l'inventaire range la cote parmi les papiers du Séminaire
de géométrie algébrique, et l'on s'en tient là.}

\pagerange{3}{3}
\section*{Sections monogènes et le foncteur qu'elles forment (page 3)}

Soit $f : X \to Y$ de type fini. On dit que $X$ est \emph{monogène} sur $Y$ si
ses fibres sont géométriquement irréductibles et si $f$ est ouvert ; qu'une
application $\varphi : Y \to X$ est \emph{monogène} s'il existe un ouvert $U$
de $X$, universellement ouvert sur $Y$, tel que pour tout $y$, $U \cap X_y$
soit non vide et géométriquement irréductible, de point générique
$\varphi(y)$.\footnote{La page écrit d'abord seulement « ouvert », puis reprend
la ligne et donne « universellement ouvert sur $Y$ », qui est aussi ce que
retient la définition 3 de la page 11. La non-vacuité et l'irréductibilité
géométrique de $U \cap X_y$ ne sont pas écrites dans la définition de la page
3, mais sans elles « le point générique de $U \cap f^{-1}(y)$ » n'a pas de sens
et la stabilité par changement de base qui suit est fausse ; on les ajoute.
Une variante « si $Y$ est noethérien, il revient au même de dire » suit, dont
la seconde condition est illisible, ainsi qu'un bloc marginal en biais qui ne
se lit pas au-delà de l'implication qui l'ouvre.}

\emph{Stabilité par changement de base.} Si $\varphi$ est monogène, et
$Y' \to Y$ un morphisme quelconque, posons $X' = X \times_Y Y'$ et
$U' = U \times_Y Y'$. Alors $U'$ est ouvert dans $X'$ et encore universellement
ouvert sur $Y'$ ; la fibre de $U'$ en $y'$ au-dessus de $y$ est
$U_y \otimes_{\kappa(y)} \kappa(y')$, irréductible parce que $U_y$ est
géométriquement irréductible. On définit $\varphi'(y')$ comme son point
générique, et $\varphi'$ est monogène.\footnote{Ce n'est pas la composée de
$\varphi$ et de la projection --- un point générique ne se transporte pas par
image réciproque ---, d'où son « définie de façon évidente ». C'est là que
l'irréductibilité \emph{géométrique} des fibres, et non leur seule
irréductibilité, est indispensable.} Ainsi
\[
  F(Y') = \{\text{sections monogènes de } X' / Y'\}
\]
est un foncteur contravariant en $Y'$.

\emph{La question.} $F$ est-il représentable ? Si oui, remarque-t-il, le
représentant $M$ est vraisemblablement étale de type fini sur $Y$, mais pas
nécessairement séparé.

\emph{Le critère.} Soit $M$ étale de type fini sur $Y$, et $\varphi_M$ une
application monogène pour $X \times_Y M / M$. Supposons que pour tout $y \in Y$
et toute extension algébriquement close $\Omega$ de $\kappa(y)$, l'application
qui à un point $m \in M(\Omega)$ associe la composante de $X \otimes_Y \Omega$
de point générique $\varphi_M(m)$,
\[
  M(\Omega) \longrightarrow \mathrm{Irr}(X \otimes_Y \Omega),
\]
soit bijective. Alors $(M, \varphi_M)$ représente $F$. La page ne le démontre
pas ; la page 4 dit à quoi la démonstration se ramène.\footnote{Une note en
marge le glose : « $M$ correspond aux points génériques des fibres de
$X \to Y$ ». Le nom moderne de l'objet est \emph{espace des composantes
irréductibles} de $X/Y$ ; il est traité dans M.~Romagny, « Composantes
connexes et irréductibles en familles », \emph{Manuscripta Math.} 136 (2011),
dans le cadre des espaces algébriques --- ce que la remarque « pas
nécessairement séparé » laisse prévoir. Le dossier ne cite personne, et la
parenté n'est pas une filiation.}

\pagerange{4}{4}
\section*{Réduction au cas d'une base locale complète (page 4)}

\emph{Ce qu'il faut prouver.} Soit $\psi : Y \to M$ une application (d'ensembles)
telle que $\psi(y)$ soit rationnel sur $\kappa(y)$ pour tout $y$, et que
l'application $\varphi : Y \to X$ qui à $y$ associe le point générique de la
composante géométriquement irréductible de $X_y$ indexée par $\psi(y)$ soit
monogène. Alors $\psi(Y)$ est ouvert dans $M$, et $\psi$ provient d'une section
de $M$ sur $Y$.\footnote{La page écrit $\varphi : Y' \to X'$ dans un énoncé où
tout se passe sur $Y$ ; on lit $Y$. Elle dit la conclusion « évidente, grâce à
la condition de rationalité ». La seconde moitié l'est en effet : un ouvert
$V = \psi(Y)$ de $M$ qui s'envoie bijectivement sur $Y$ avec
$\kappa(\psi(y)) = \kappa(y)$ est étale et radiciel et surjectif sur $Y$, donc
isomorphe à $Y$. Que $\psi(Y)$ soit ouvert, c'est ce que la monogénéité doit
fournir, et la page ne dit pas comment.}

\emph{Conséquence : le problème est local.} Si le problème des modules a une
solution au-dessus du schéma local $\mathrm{Spec}(\mathcal{O}_{Y,y})$, il en a
une au-dessus d'un voisinage ouvert de $y$.\footnote{La page écrit
$\mathrm{Spec}(\Omega_y)$, deux fois sur la ligne. Au-dessus d'un point
géométrique la solution est tautologique --- l'ensemble des composantes de la
fibre géométrique --- et l'énoncé ne dirait rien ; la phrase suivante de la
page, « cela nous ramène au cas où $Y = \mathrm{Spec}(A)$, $A$ local », montre
que c'est le schéma local qui est en jeu. On le lit ainsi. La page n'argumente
pas le passage du schéma local à un voisinage ; c'est un argument de passage à
la limite pour un $M$ de présentation finie, mais la propriété de
représentation elle-même doit encore être propagée aux points voisins, et rien
ici ne le fait.} On se ramène ainsi, pour l'existence, au cas
$Y = \mathrm{Spec}(A)$ avec $A$ local.

\emph{Passage au complété.} Par descente des morphismes étales, on se ramène
ensuite au cas où $A$ est complet, et l'on peut supposer que la fibre fermée se
décompose en composantes géométriquement irréductibles.\footnote{La descente
le long de $\mathrm{Spec}(\widehat{A}) \to \mathrm{Spec}(A)$, fidèlement plat,
est effective pour un $M$ séparé : un morphisme quasi-fini séparé est
quasi-affine (théorème principal de Zariski), et les données de descente sur un
schéma quasi-affine sont effectives. Pour un $M$ non séparé, que la page 3
vient d'envisager, l'effectivité n'est pas garantie dans la catégorie des
schémas ; c'est l'une des raisons pour lesquelles la théorie moderne travaille
avec des espaces algébriques. La page ne fait pas la distinction. Qu'on puisse
rendre les composantes de la fibre fermée géométriquement irréductibles, c'est
une extension finie du corps résiduel, relevée à $A$ complet par un
revêtement étale ; la page l'affirme sans le détailler, et plusieurs mots de la
ligne sont illisibles.}

\pagerange{5}{6}
\section*{La construction quand la fibre fermée est irréductible (pages 5 et 6)}

Hypothèses : $f : X \to Y = \mathrm{Spec}(A)$ universellement ouvert, de type
fini, $A$ local noethérien complet, $y$ le point fermé, et $X_y$
irréductible.\footnote{Un adjectif est ajouté devant « universellement
ouvert », lu « strictement » avec doute et glosé dans un bloc marginal en biais
qui ne se lit pas ; on ne le reprend pas. Un « N.B. » entouré surmonte « de
type fini », sans texte lisible.}

\emph{Premier temps : exclure la ramification.} Soit $x$ un point de la fibre
fermée. La page commence par supposer que $x$ est dans l'adhérence de deux
composantes distinctes d'une fibre, et construit, pour réfuter cette
éventualité, un trait $Y' = \mathrm{Spec}(W)$ au-dessus de $Y$, $W$ anneau de
valuation discrète, en normalisant $A$ dans une extension finie $K'$ du corps
des fractions $K$ : l'anneau $A'$ obtenu est fini sur $A$ et \emph{local}, parce
que $A$ est local complet, donc hensélien.\footnote{Il écrit « $\frac{1}{2}$
local complet », le $\frac{1}{2}$ étant son abréviation de « semi- » : un
anneau fini sur un anneau local complet est semi-local complet, donc produit de
locaux, et $A'$ intègre est local. On restitue l'argument dans ce sens. Le
premier alinéa (i) est presque entièrement illisible, et le « lemme b) » qu'il
invoque n'est pas identifié dans le dossier.} Le cas est déclaré « absurde ».

\emph{Second temps : la section.} Pour chaque $y \in Y$ il y a donc une unique
composante $\xi \in E_{X/Y}(y)$ dont l'adhérence contient $x$, et elle est
géométriquement irréductible. On pose $\varphi(y) = $ son point générique. Les
deux conditions d'une application monogène sont satisfaites, et il existe un
ouvert $U$ de $X$ contenant $\varphi(Y)$ dont les fibres sur $Y$ sont
géométriquement irréductibles.

\emph{Reste la continuité} (page 6), que la page ramène à l'ouverture de
$X \to Y$, et c'est exactement ce que la Prop.~1 de la page 11 établira :
pour la section des points génériques d'un morphisme à fibres irréductibles,
continuité et ouverture sont une seule et même chose.\footnote{Un bloc de
deux lignes enfermé dans une boucle est annulé au milieu de la page 6. La page
enchaîne sur le cas où la fibre fermée n'est plus irréductible : on traite
chacune de ses composantes $i = 1, \ldots, n$ par la construction précédente,
et l'on recolle les $n$ sections obtenues dans un $M'$ ; les mots qui disent
comment sont presque tous illisibles, et l'on n'en tire pas davantage.}

\pagerange{7}{7}
\section*{Le Lemme de la page 7}

La page récapitule sous forme de lemme ce qui vient d'être fait. Soit $X$ de
type fini sur $Y = \mathrm{Spec}(A)$, avec $A$ local complet et $X$
universellement ouvert sur $Y$, et supposons que les composantes de la fibre
fermée soient géométriquement irréductibles. Il existe alors un $Y$-schéma $M$
qui représente les composantes géométriquement irréductibles des fibres de
$X/Y$ ; au-dessus de l'ouvert $Y' = Y - \{y\}$ il est séparé et
fini.\footnote{Des cinq hypothèses du Lemme, seules (i), (ii) et (iii) se
lisent ; (iv) parle de la disparition de composantes au-dessus de $Y'$ et (v)
de recollements des composantes suivis d'une extension finie de la base, mais
la plupart de leurs mots sont illisibles, et l'on ne les restitue pas. Le nom
du $Y$-schéma est d'une terminaison incertaine.}

La page ajoute un critère de séparation, sous un titre souligné que la
transcription lit « Le Moyen » sans l'assurer : si, après tout changement de
base local $A \to A_1$, aucune composante d'une fibre de $X \times_Y Y_1$ ne se
spécialise en deux composantes distinctes, alors $M/Y$ est
séparé.\footnote{L'énoncé est fait de fragments séparés par des mots
illisibles ; on en donne le sens que ses mots lisibles --- « se spécialise
en », « deux distinctes », « Alors $M/Y$ est séparé » --- permettent, et qui
est celui de la condition (ii) de la page 8.}

\pagerange{8}{9}
\section*{Le théorème (pages 8 et 9)}

\emph{L'énoncé du feuillet.} Soit $f : X \to Y$ de type fini. On suppose que
pour tout point générique $x$ d'une composante d'une fibre $X_y$, il existe un
voisinage de $x$ qui soit universellement ouvert et quasi-fini sur
$Y[t_1, \ldots, t_{n-1}]$ --- condition vérifiée, dit-il, si $X$ est plat sur
$Y$ au voisinage de $x$.\footnote{« Principal » est son mot pour ces points
génériques. L'indice $n - 1$ est sur le feuillet sans que $n$ soit défini ; on
comprend la dimension de la fibre en $x$ augmentée de un, mais c'est une
lecture. L'assertion sur le cas plat n'est pas démontrée ici et cette lecture ne
la vérifie pas : ce qu'on sait sans peine est qu'un morphisme plat de
présentation finie est universellement ouvert, et c'est cette propriété que
la construction des pages 5 et 6 utilise.} Pour qu'il existe un $M$
\emph{étale et séparé} sur $Y$ indexant les composantes géométriquement
irréductibles des fibres de $X/Y$, il faut et il suffit que, pour tout trait
$Y' = \mathrm{Spec}(W) \to Y$, de point générique $y'$ et de point fermé
$\bar{y}'$, en posant $X' = X \times_Y Y'$ :

\begin{enumerate}
\item[(ii)] aucune composante de la fibre générique $X'_{y'}$ ne se spécialise
en deux composantes irréductibles distinctes de la fibre spéciale ;
\item[(iii)] deux composantes irréductibles distinctes ne se rencontrent pas
dans une même spécialisation.
\end{enumerate}

La démonstration n'est pas sur les feuillets, et la première condition du
théorème, qui précède (ii), est illisible.\footnote{La condition (ii) est
donnée avec la glose qu'il écrit en interligne, « i.e.\ une comp.\ de la fibre
générique ne se spécialise en deux composantes abs.\ irréductibles distinctes
de la fibre spéciale ». La condition (iii) est lacunaire : on lit
« composantes abs.\ irréductibles distinctes », « ne », « spécialiser en »,
« comp.\ abs.\ irréd.\ de $f^{-1}(y')$ » ; le sens qu'on lui donne ici est
celui que ces mots permettent et que l'analyse qui suit rend nécessaire, et il
est le nôtre. Les hypothèses sur le trait --- « fini et plat », « $V^{\circ}$
fini » --- ne se lisent qu'en fragments.}

\emph{Pourquoi ces conditions : les trois critères valuatifs.} Le théorème a la
forme qu'il a parce qu'un schéma étale au-dessus d'un trait se lit sur ses
points. Supposons $W$ strictement hensélien et $M$ étale de type fini sur
$Y' = \mathrm{Spec}(W)$ ; ce qui suit est notre analyse, non celle du
feuillet.

\begin{itemize}
\item \emph{Toute étale.} Si $m$ est un point de la fibre spéciale de $M$,
l'anneau local $\mathcal{O}_{M,m}$ est une $W$-algèbre locale étale de même
corps résiduel, donc égale à $W$ : un seul point générique de $M$ se spécialise
en $m$. Dans le langage des composantes, une composante de la fibre spéciale
provient d'une seule composante générique --- deux composantes génériques ne
viennent pas se fondre en une.
\item \emph{Séparation.} Par le critère valuatif, $M$ est séparé si et seulement
si un point générique se spécialise en au plus un point spécial : une
composante générique ne se casse pas en deux. C'est la condition (ii).
\item \emph{Propreté.} Par le critère valuatif d'existence, $M$ (séparé, de
type fini) est propre si et seulement si tout point générique se spécialise en
au moins un point spécial : aucune composante générique ne disparaît. Un
schéma étale, séparé et propre étant fini, $M$ est alors un revêtement étale
fini de $Y$.
\end{itemize}

\emph{Le N.B.\ et le corollaire} (page 9). Sous les conditions du théorème,
l'ensemble des composantes de la fibre spéciale s'injecte dans celui des
composantes de la fibre générique --- à chaque point spécial de $M$, le Lemme de
Hensel attache une unique section de $M$ sur le trait, et le point générique de
cette section ---, mais l'injection n'est pas nécessairement
surjective.\footnote{La page écrit que les composantes de $f^{-1}(\bar{y}')$
s'injectent dans celles de $f^{-1}(y')$ ; on lit la barre comme désignant le
point fermé du trait, ce que la suite de la page, où $y'$ est générique et $y$
spécial, impose. La justification par les sections de $M$ est la nôtre.}
D'où le complément : pour que $M$ soit en outre fini sur $Y$, il faut et il
suffit que

\begin{enumerate}
\item[(iii)] pour tout point générique $\xi'$ d'une composante de la fibre
générique, il existe un point générique $x'$ d'une composante de la fibre
spéciale tel que $x' \in \overline{\{\xi'\}}$.
\end{enumerate}

C'est exactement le critère d'existence ci-dessus.\footnote{La page annonce ce
complément sous un mot souligné lu « Cor. » sans certitude, et l'énonce comme la
condition pour que $M$ soit « un \emph{ouvert} étale sur $Y$ ». Ce que la
condition fournit est la propreté, donc la finitude, et non une ouverture ; le
mot, tel qu'il se lit, ne porte pas le sens de la condition, et l'on énonce ce
que la condition donne. Il numérote ce complément (iii), comme la seconde
condition de la page 8 ; on garde ses numéros.}

\pagerange{10}{10}
\section*{Deux remarques (page 10)}

\begin{enumerate}
\item[1)] La première condition du théorème, seule, est nécessaire et suffisante
pour l'existence d'un $M$ étale \emph{non nécessairement séparé}. C'est
l'analogue, dans l'analyse de la page précédente, du premier des trois critères
: ce qu'il faut pour que $M$ soit étale, avant toute question de séparation ou
de propreté.\footnote{Le rapprochement est le nôtre, et il repose sur une
condition que la page 8 ne laisse pas lire. Au-dessus d'un mot « Corollaire »
biffé.}
\item[2)] La condition (iii) est satisfaite si $f$ est séparable ; et si $f$
est propre, $M$ est propre sur $Y$.\footnote{Entre les deux, la page ramène la
conjonction de (ii) et (iii), dans le cas propre, à trois conditions a), b),
c), dont on ne lit que des bribes (« d'une composante rel.\ irréd. », « une
spécialisation », « $M$ existe ») ; on ne les restitue pas. La seconde
assertion est cohérente avec l'analyse de la page 9 : si $f$ est propre,
l'adhérence d'une composante générique rencontre la fibre spéciale, et, $f$
étant de plus universellement ouvert, contient une composante de celle-ci.
L'argument est le nôtre.}
\end{enumerate}

\pagerange{11}{11}
\section*{Fibres irréductibles : continuité et ouverture (page 11)}

\textbf{Proposition 1.} Soit $f : X \to Y$ un morphisme de type fini dont les
fibres sont irréductibles et non vides, et soit $\varphi(y)$ le point générique
de $X_y$. On considère les conditions :

\begin{enumerate}
\item[(i)] $\varphi : Y \to X$ est continue ;
\item[(ii)] $f$ est ouvert ;
\item[(iii)] $y \in \overline{\{y'\}} \Longrightarrow \varphi(y) \in \overline{\{\varphi(y')\}}$ ;
\item[(iii bis)] pour tous $y \in \overline{\{y'\}}$ et $x$ au-dessus de $y$, il
existe $x'$ au-dessus de $y'$ tel que $x \in \overline{\{x'\}}$.
\end{enumerate}

Alors (i) $\Leftrightarrow$ (ii) sans autre hypothèse, et
(ii) $\Rightarrow$ (iii) $\Leftrightarrow$ (iii bis) ; si $f$ est localement de
présentation finie --- par exemple $Y$ localement noethérien ---, les quatre
conditions sont équivalentes.

\emph{Démonstration.} Pour $U$ ouvert dans $X$, le point générique de la fibre
irréductible $X_y$ est dans $U$ si et seulement si $U$ rencontre $X_y$ ; donc
\[
  \varphi^{-1}(U) = f(U),
\]
et $\varphi$ est continue si et seulement si $f$ est ouvert. Une application
continue respecte la spécialisation, d'où (i) $\Rightarrow$ (iii). Si (iii)
vaut, et $x$ est au-dessus de $y \in \overline{\{y'\}}$, on a
$x \in \overline{\{\varphi(y)\}} \subset \overline{\{\varphi(y')\}}$, et
$x' = \varphi(y')$ convient : c'est (iii bis). Inversement (iii bis) appliquée
à $x = \varphi(y)$ donne $\varphi(y) \in \overline{\{x'\}}$ avec $x' \in X_{y'}$,
donc $x' \in \overline{\{\varphi(y')\}}$, d'où (iii). Enfin (iii bis) dit que les
générisations se relèvent le long de $f$, et pour un morphisme localement de
présentation finie cela entraîne que $f$ est ouvert.\footnote{La formule
$\varphi^{-1}(U) = f(U)$ est sur la page ; le reste de la démonstration est
presque entièrement illisible, et on la rédige. La page énonce (ii) sous la
forme « $f$ soit ouvert, […] si $Y$ est loc.\ noeth., $f$ de type fini », ce qui
place l'hypothèse de finitude au bon endroit : elle ne sert qu'au retour de
(iii bis) vers (ii), qui est le relèvement des générisations
(EGA~IV, 1.10.4). La condition (iii bis) est encadrée sur le feuillet.}

\textbf{Définition 2.} $f$ est \emph{universellement monogène} s'il est
universellement ouvert et à fibres géométriquement irréductibles --- de sorte
que la Proposition 1 s'applique après tout changement de
base.\footnote{La page n'écrit que « Définition de \emph{universellement
monogène} » après une première phrase inachevée ; on complète par ce que la
page 3 et la page 12 en utilisent.}

\textbf{Définition 3.} Une application $\varphi : Y \to X$ est \emph{monogène}
s'il existe un ouvert $U$ de $X$, universellement monogène sur $Y$, dont
$\varphi$ est la section des points génériques.

\pagerange{12}{12}
\section*{La question reposée (page 12)}

La propriété monogène est locale sur $Y$ : elle dit que $\varphi(y)$ est le
point générique d'une composante de $X_y$, et qu'au voisinage il existe un
ouvert $U$ universellement ouvert dont les fibres sont géométriquement
irréductibles.\footnote{La page ajoute qu'on peut rétrécir $U$ en un
$U' \subset U$ encore surjectif sur $Y$ --- lecture incertaine ---, et plusieurs
mots de la phrase sont illisibles.}

Notant $I(X/S)$ l'ensemble des sections monogènes de $X$ sur $S$, le foncteur
$S' \mapsto I(X \times_S S' / S')$ est contravariant de façon évidente. « Est-il
représentable ??? » La question de la page 3 revient donc, cette fois sur une
base notée $S$, et le dossier ne la tranche pas.\footnote{Dans l'état des
feuillets, la représentabilité n'est acquise que sous les conditions du
théorème de la page 8, qui n'est pas démontré. La réponse moderne est qu'en
général le foncteur n'est pas représentable par un schéma, mais qu'il l'est
par un espace algébrique étale sous des hypothèses de platitude et de
présentation finie ; on renvoie à l'article de Romagny cité plus haut pour
l'énoncé exact, que cette lecture ne reproduit pas.}

La page se ferme sur un cube de changements de base, au crayon : $X \to S$,
son changement de base $X_M \to M$ par $M \to S$, et la face obtenue au-dessus
d'un point géométrique $S' = \mathrm{Spec}(\Omega)$, avec $X'_{m'} \to m'$.

\begin{tikzcd}[column sep=small, row sep=small, nodes={font=\scriptsize}]
X \arrow[dd] & & X_M \arrow[ll] \arrow[dd] & \\
& X' \arrow[ul] \arrow[dd] & & X'_{m'} \arrow[ll] \arrow[ul] \arrow[dd] \\
S & & M \arrow[ll] & \\
& S' = \mathrm{Spec}(\Omega) \arrow[ul] & & m' \arrow[ll] \arrow[ul]
\end{tikzcd}

C'est la situation du critère de la page 3 : le point $m'$ de $M$ au-dessus du
point géométrique $S'$ doit indexer une composante de $X'$.\footnote{Le montant
de gauche du cube porte $\kappa(s)$ en haut et $\kappa(s')$ en bas, reliés par
un trait à $S$ et à $S'$.}

\pagerange{13}{13}
\section*{Un lemme barré, et la spécialisation des composantes (page 13)}

\emph{Le lemme barré.} Soit $f : X \to Y$ plat de type fini, $Y$ local
noethérien, $y$ un point de $Y$. Il existerait un ouvert $U$ tel que, pour tout
$y'$ et tout point générique $x'$ d'une composante de $X_{y'} \cap U$, un point
générique $x$ d'une composante de $X_y$ soit spécialisation de $x'$. Il est
annulé par deux grandes diagonales, et il ne tient pas tel
quel.\footnote{Contre-exemple, si $U$ doit contenir la fibre $X_y$ tout
entière --- la condition sur $U$ n'est pas lisible : sur un trait
$\mathrm{Spec}(W)$ d'uniformisante $t$, soit $X \subset \mathbb{A}^3_W$ la
réunion du plan $z = 0$ et de la droite $x = 0$, $z = t$. Chacune des deux
composantes domine le trait et $X$ est réduit, donc $X$ est plat sur $W$. La
fibre générique a deux composantes, le plan et une droite disjointe ; dans la
fibre spéciale la droite vient se coucher dans le plan, dont elle n'est plus
une composante. Le point générique de la droite générique n'a donc aucune
composante de la fibre spéciale parmi ses spécialisations. Il faudrait une
hypothèse d'équidimensionalité.}

\emph{Le changement de base et les spécialisations.} Deux diagrammes
cartésiens, $X' = X \times_Y Y'$ au-dessus de $Y' \to Y$, puis
$X'' = X \times_Y Y''$ au-dessus de $Y'' \to Y' \to Y$, portent chacun leurs
points : $x \leftarrow \xi$ au-dessus de $y \leftarrow \eta$,
$x' \leftarrow \xi'$ au-dessus de $y' \leftarrow \eta'$.

\begin{tikzcd}[column sep=small, row sep=small, nodes={font=\scriptsize}]
X \arrow[d, "f"] & X' = X \times_Y Y' \arrow[l] \arrow[d, "f'"] & X'' = X \times_Y Y'' \arrow[l] \arrow[d] \\
Y & Y' \arrow[l] & Y'' \arrow[l]
\end{tikzcd}

\emph{Hypothèses} : $x \in E_{X/Y}(y)$, $\xi \in E_{X/Y}(\eta)$,
$x \in \overline{\{\xi\}}$ ; $y'$ au-dessus de $y$, $\eta'$ au-dessus de
$\eta$, $y' \in \overline{\{\eta'\}}$. \emph{Conclusions espérées} :

\begin{enumerate}
\item[a)] pour tout $x' \in E_{X'/Y'}(y')$ au-dessus de $x$, il existe
$\xi' \in E_{X'/Y'}(\eta')$ au-dessus de $\xi$ avec
$x' \in \overline{\{\xi'\}}$ ;
\item[b)] pour tout $\xi' \in E_{X'/Y'}(\eta')$ au-dessus de $\xi$, il existe
$x' \in E_{X'/Y'}(y')$ au-dessus de $x$ avec $x' \in \overline{\{\xi'\}}$.
\end{enumerate}

Ce sont les deux sens --- montée et descente --- de la compatibilité entre la
spécialisation des composantes et le changement de base, dont la définition de
$M$ comme foncteur a besoin. Il écrit « Vérifier ! », et ne vérifie
pas.\footnote{Une flèche venue de la marge porte une hypothèse supplémentaire
sur $Y' \to Y$ dont les mots sont illisibles. Un « Concl' » qui suivait,
formulé par deux applications
$\Gamma_{X'/Y'}(y', \eta') \to E_{X'/Y'}(y')$ et
$\Gamma_{X'/Y'}(y', \eta') \to E_{X'/Y'}(\eta')$ dont l'une au moins devait
être surjective, est annulé par une grande boucle. Sans hypothèse sur
$Y' \to Y$ (plate, par exemple) cette lecture ne garantit ni a) ni b).}

\end{document}
